\documentclass[../main]{subfiles}
\begin{document}

\section{Práctico de Mantenimiento: Motosierras}

\info{Mantenimiento de Motosierras}{La motosierra es una herramienta fundamental en tareas rurales, forestales y de jardinería. Su correcto mantenimiento no solo prolonga la vida útil del equipo, sino que garantiza la seguridad del operador.}

\subsection{Mantenimiento Preventivo Básico}

Antes de realizar cualquier ajuste en el motor, es fundamental revisar los siguientes componentes:

\begin{itemize}
    \item \textbf{Filtro de aire:} Debe estar limpio. Un filtro obstruido altera la mezcla de aire y combustible, provocando un funcionamiento irregular.
    \item \textbf{Bujía:} Revisar el estado del electrodo y la separación. Si está sucia de carbón, debe limpiarse o reemplazarse para asegurar una chispa fuerte.
    \item \textbf{Filtro de combustible:} Se encuentra en el tanque y debe cambiarse periódicamente para evitar obstrucciones en el carburador.
    \item \textbf{Lubricación de cadena:} Es vital verificar que el aceite llegue correctamente a la cadena y la espada para evitar el sobrecalentamiento y el desgaste prematuro.
    \item \textbf{Afiliado de la cadena:} Una cadena mal afilada obliga al motor a trabajar en exceso, generando sobrecalentamiento y pérdida de rendimiento.
\end{itemize}

\subsection{Regulación del Carburador: Tornillos H y L}

El carburador tiene la función de mezclar el aire y el combustible en la proporción correcta. Para ajustar esta mezcla, la mayoría de las motosierras cuentan con tres tornillos de regulación:

\begin{itemize}
    \item \textbf{Tornillo L (Low / Bajas):} Regula la mezcla de aire y combustible a bajas revoluciones (ralentí). Controla la respuesta del motor al iniciar la aceleración.
    \item \textbf{Tornillo H (High / Altas):} Controla la mezcla a altas revoluciones (potencia máxima). Determina la velocidad máxima del motor bajo carga.
    \item \textbf{Tornillo T o LA (Ralentí):} Ajusta la velocidad del motor cuando el acelerador está suelto, evitando que la cadena se mueva o que el motor se apague.
\end{itemize}

\subsubsection{Pasos para el ajuste del carburador}

Antes de comenzar, asegúrese de que el motor esté caliente tras unos minutos de funcionamiento y que el filtro de aire esté limpio.

\begin{enumerate}
    \item \textbf{Ajuste inicial:} Gire suavemente los tornillos H y L en sentido horario hasta que hagan tope. \textbf{No los apriete con fuerza} para no dañar los asientos del carburador. Luego, gírelos en sentido antihorario \textbf{1 vuelta y media} como punto de partida (o según el manual del fabricante).
    \item \textbf{Ajuste del tornillo L (Bajas):} Con la motosierra encendida en ralentí, gire el tornillo L lentamente en sentido horario hasta que el motor acelere y luego comience a ahogarse. En ese punto, gire en sentido antihorario hasta que el motor funcione de manera suave y uniforme.
    \item \textbf{Ajuste del tornillo T (Ralentí):} Ajuste este tornillo para que la cadena no gire cuando el motor está en reposo, pero evitando que el motor se apague.
    \item \textbf{Ajuste del tornillo H (Altas):} Acelere la motosierra a fondo. Si el motor suena agudo y chillón, la mezcla es pobre (girar H en sentido antihorario para dar más gas). Si suena grave y humea excesivamente, la mezcla es rica (girar H en sentido horario para dar más aire). El ajuste ideal debe permitir una aceleración rápida y un sonido limpio.
\end{enumerate}

\info{Advertencia de Seguridad}{Un ajuste incorrecto del tornillo H (mezcla demasiado pobre) puede llevar a un sobrecalentamiento severo y fundir el pistón por falta de lubricación. Si no está seguro, es preferible dejar la mezcla ligeramente rica.}

\newpage

\actividad{Leer, interpretar y contestar}
\begin{enumerate}
    \item ¿Cuál es la función específica de los tornillos H, L y T en el carburador de una motosierra?
    \item ¿Por qué es fundamental revisar el filtro de aire y la bujía antes de intentar ajustar el carburador?
    \item ¿Qué riesgos conlleva una regulación incorrecta del tornillo H?
    \item Describa brevemente el procedimiento para ajustar el tornillo L (bajas revoluciones).
\end{enumerate}

\end{document}